\section{Key Results}
% Sources: notebook cells 16, 18, 20, 24 (cross-validation), 26 and 29 (test set).
% Final-model test values are from classification_report; unweighted test values from its confusion matrix (0 of 50 strokes).
Table~\ref{tab:results} compares the experiments with 5-fold cross-validation on the training rows, followed by the single test-set evaluation.
Accuracy is shown but not used to choose the model, because the unweighted model matches the 95.13\% majority baseline while finding almost no strokes.
All metrics come from the classification confusion matrix: recall is the share of real strokes the model catches, precision is the share of flagged patients who really had a stroke, and F1 is their harmonic mean.
F1 runs from 0 to 1 (best) and is only high when both are high. ROC-AUC is threshold-free: it measures how well the predicted probabilities rank stroke patients above non-stroke patients.

\begin{table}[!htb]
\centering
\begin{tabular}{lcccc}
\hline
Model & Accuracy & Precision & Recall & F1 \\
\hline
\multicolumn{5}{l}{\textit{5-fold cross-validation, training rows (4,087 patients)}} \\
Majority baseline (always `no stroke') & 0.951 & -- & 0.000 & 0.000 \\
Baseline LogReg, unweighted ($C=1$) & 0.951 & 0.200 & 0.005 & 0.010 \\
LogReg, balanced ($C=1$) & 0.739 & 0.133 & 0.789 & 0.228 \\
LogReg, balanced, tuned ($C=0.01$) & 0.737 & 0.137 & 0.824 & 0.234 \\
\hline
\multicolumn{5}{l}{\textit{Held-out test set (1,022 patients, 50 strokes), threshold 0.5}} \\
Baseline LogReg, unweighted ($C=1$) & 0.951 & 0.000 & 0.000 & 0.000 \\
Final model (balanced, $C=0.01$) & 0.76 & 0.13 & 0.72 & 0.22 \\
\hline
\end{tabular}
\caption{Performance of each experiment. LogReg - logistic regression classifier. The majority baseline never predicts a stroke, so its precision is undefined (0/0) and its recall and F1 are zero by definition.}
\label{tab:results}
\end{table}

Class weighting made the biggest difference: cross-validated recall rose from 0.005 to 0.789 and F1 from 0.010 to 0.228, while tuning $C$ added only a little (F1 0.234).

\textbf{Threshold.}
The decision threshold was chosen from cross-validated probabilities on the training rows (Figure~\ref{fig:thresh}).
Lowering the threshold raises recall but flags more patients for a follow-up.
A threshold of 0.5 flags the fewest patients (29.5\%) while keeping recall above 0.8 (0.824); raising it to 0.6 would cut recall to 0.683 with only a small gain in precision (0.136 to 0.160).

% Two plots side by side: each minipage gets its own caption and label.
\begin{figure}[!htb]
\centering
\begin{minipage}[b]{0.41\textwidth}
\centering
\includegraphics[width=\textwidth]{threshold.png}
\caption{Cross-validated precision, recall and share flagged against the decision threshold (training rows).}
\label{fig:thresh}
\end{minipage}\hfill
\begin{minipage}[b]{0.56\textwidth}
\centering
\includegraphics[width=\textwidth]{confusion_matrix.png}
\caption{Test-set confusion matrices for the unweighted model and the final balanced model.}
\label{fig:cm}
\end{minipage}
\end{figure}

\textbf{Test set.}
On the test set the unweighted model missed all 50 strokes, whereas the final model caught 36 of them (Figure~\ref{fig:cm}), with a ROC-AUC of 0.826.
The cost is 235 false alarms, so about 1 in 7.5 flagged patients had a stroke; in screening, a missed stroke is far worse than an extra check, so this trade-off is acceptable.

\textbf{Efficiency.}
The hardware used for this exercise is a Lenovo Legion ACH6H (AMD Ryzen 7 5800H) with 32GB RAM, running the Linux distro Pop!\_OS with Python 3.12.3 and sklearn 1.9.1; timings vary between runs and with background load, so they are approximate.
Fitting the final model on 4,087 rows took about 19\,ms, predicting 1,022 patients took about 5\,ms, and peak Python memory during fitting (measured with \texttt{tracemalloc}, which does not count all C-level allocations) was 1.61\,MB.
Balanced weighting made no meaningful difference to cost (unweighted: 21.8\,ms fit, 5.4\,ms predict, 1.61\,MB).
Fit time grew roughly linearly with training size, from about 13\,ms at 650 rows to 18.5\,ms at 3,270 rows (Figure~\ref{fig:lc}), but this is only 1.4 times the cost for 5 times the rows, so a fixed overhead of about 11\,ms dominates and the scaling evidence is limited.
